% Options for packages loaded elsewhere
% Options for packages loaded elsewhere
\PassOptionsToPackage{unicode}{hyperref}
\PassOptionsToPackage{hyphens}{url}
\PassOptionsToPackage{dvipsnames,svgnames,x11names}{xcolor}
%
\documentclass[
  letterpaper,
  DIV=11,
  numbers=noendperiod]{scrartcl}
\usepackage{xcolor}
\usepackage{amsmath,amssymb}
\setcounter{secnumdepth}{-\maxdimen} % remove section numbering
\usepackage{iftex}
\ifPDFTeX
  \usepackage[T1]{fontenc}
  \usepackage[utf8]{inputenc}
  \usepackage{textcomp} % provide euro and other symbols
\else % if luatex or xetex
  \usepackage{unicode-math} % this also loads fontspec
  \defaultfontfeatures{Scale=MatchLowercase}
  \defaultfontfeatures[\rmfamily]{Ligatures=TeX,Scale=1}
\fi
\usepackage{lmodern}
\ifPDFTeX\else
  % xetex/luatex font selection
\fi
% Use upquote if available, for straight quotes in verbatim environments
\IfFileExists{upquote.sty}{\usepackage{upquote}}{}
\IfFileExists{microtype.sty}{% use microtype if available
  \usepackage[]{microtype}
  \UseMicrotypeSet[protrusion]{basicmath} % disable protrusion for tt fonts
}{}
\makeatletter
\@ifundefined{KOMAClassName}{% if non-KOMA class
  \IfFileExists{parskip.sty}{%
    \usepackage{parskip}
  }{% else
    \setlength{\parindent}{0pt}
    \setlength{\parskip}{6pt plus 2pt minus 1pt}}
}{% if KOMA class
  \KOMAoptions{parskip=half}}
\makeatother
% Make \paragraph and \subparagraph free-standing
\makeatletter
\ifx\paragraph\undefined\else
  \let\oldparagraph\paragraph
  \renewcommand{\paragraph}{
    \@ifstar
      \xxxParagraphStar
      \xxxParagraphNoStar
  }
  \newcommand{\xxxParagraphStar}[1]{\oldparagraph*{#1}\mbox{}}
  \newcommand{\xxxParagraphNoStar}[1]{\oldparagraph{#1}\mbox{}}
\fi
\ifx\subparagraph\undefined\else
  \let\oldsubparagraph\subparagraph
  \renewcommand{\subparagraph}{
    \@ifstar
      \xxxSubParagraphStar
      \xxxSubParagraphNoStar
  }
  \newcommand{\xxxSubParagraphStar}[1]{\oldsubparagraph*{#1}\mbox{}}
  \newcommand{\xxxSubParagraphNoStar}[1]{\oldsubparagraph{#1}\mbox{}}
\fi
\makeatother


\usepackage{longtable,booktabs,array}
\newcounter{none} % for unnumbered tables
\usepackage{calc} % for calculating minipage widths
% Correct order of tables after \paragraph or \subparagraph
\usepackage{etoolbox}
\makeatletter
\patchcmd\longtable{\par}{\if@noskipsec\mbox{}\fi\par}{}{}
\makeatother
% Allow footnotes in longtable head/foot
\IfFileExists{footnotehyper.sty}{\usepackage{footnotehyper}}{\usepackage{footnote}}
\makesavenoteenv{longtable}
\usepackage{graphicx}
\makeatletter
\newsavebox\pandoc@box
\newcommand*\pandocbounded[1]{% scales image to fit in text height/width
  \sbox\pandoc@box{#1}%
  \Gscale@div\@tempa{\textheight}{\dimexpr\ht\pandoc@box+\dp\pandoc@box\relax}%
  \Gscale@div\@tempb{\linewidth}{\wd\pandoc@box}%
  \ifdim\@tempb\p@<\@tempa\p@\let\@tempa\@tempb\fi% select the smaller of both
  \ifdim\@tempa\p@<\p@\scalebox{\@tempa}{\usebox\pandoc@box}%
  \else\usebox{\pandoc@box}%
  \fi%
}
% Set default figure placement to htbp
\def\fps@figure{htbp}
\makeatother


% definitions for citeproc citations
\NewDocumentCommand\citeproctext{}{}
\NewDocumentCommand\citeproc{mm}{%
  \begingroup\def\citeproctext{#2}\cite{#1}\endgroup}
\makeatletter
 % allow citations to break across lines
 \let\@cite@ofmt\@firstofone
 % avoid brackets around text for \cite:
 \def\@biblabel#1{}
 \def\@cite#1#2{{#1\if@tempswa , #2\fi}}
\makeatother
\newlength{\cslhangindent}
\setlength{\cslhangindent}{1.5em}
\newlength{\csllabelwidth}
\setlength{\csllabelwidth}{3em}
\newenvironment{CSLReferences}[2] % #1 hanging-indent, #2 entry-spacing
 {\begin{list}{}{%
  \setlength{\itemindent}{0pt}
  \setlength{\leftmargin}{0pt}
  \setlength{\parsep}{0pt}
  % turn on hanging indent if param 1 is 1
  \ifodd #1
   \setlength{\leftmargin}{\cslhangindent}
   \setlength{\itemindent}{-1\cslhangindent}
  \fi
  % set entry spacing
  \setlength{\itemsep}{#2\baselineskip}}}
 {\end{list}}
\usepackage{calc}
\newcommand{\CSLBlock}[1]{\hfill\break\parbox[t]{\linewidth}{\strut\ignorespaces#1\strut}}
\newcommand{\CSLLeftMargin}[1]{\parbox[t]{\csllabelwidth}{\strut#1\strut}}
\newcommand{\CSLRightInline}[1]{\parbox[t]{\linewidth - \csllabelwidth}{\strut#1\strut}}
\newcommand{\CSLIndent}[1]{\hspace{\cslhangindent}#1}



\setlength{\emergencystretch}{3em} % prevent overfull lines

\providecommand{\tightlist}{%
  \setlength{\itemsep}{0pt}\setlength{\parskip}{0pt}}



 


\KOMAoption{captions}{tableheading}
\makeatletter
\@ifpackageloaded{caption}{}{\usepackage{caption}}
\AtBeginDocument{%
\ifdefined\contentsname
  \renewcommand*\contentsname{Table of contents}
\else
  \newcommand\contentsname{Table of contents}
\fi
\ifdefined\listfigurename
  \renewcommand*\listfigurename{List of Figures}
\else
  \newcommand\listfigurename{List of Figures}
\fi
\ifdefined\listtablename
  \renewcommand*\listtablename{List of Tables}
\else
  \newcommand\listtablename{List of Tables}
\fi
\ifdefined\figurename
  \renewcommand*\figurename{Figure}
\else
  \newcommand\figurename{Figure}
\fi
\ifdefined\tablename
  \renewcommand*\tablename{Table}
\else
  \newcommand\tablename{Table}
\fi
}
\@ifpackageloaded{float}{}{\usepackage{float}}
\floatstyle{ruled}
\@ifundefined{c@chapter}{\newfloat{codelisting}{h}{lop}}{\newfloat{codelisting}{h}{lop}[chapter]}
\floatname{codelisting}{Listing}
\newcommand*\listoflistings{\listof{codelisting}{List of Listings}}
\makeatother
\makeatletter
\makeatother
\makeatletter
\@ifpackageloaded{caption}{}{\usepackage{caption}}
\@ifpackageloaded{subcaption}{}{\usepackage{subcaption}}
\makeatother
\usepackage{bookmark}
\IfFileExists{xurl.sty}{\usepackage{xurl}}{} % add URL line breaks if available
\urlstyle{same}
\hypersetup{
  pdftitle={Klippel-Trenaunay Syndrome Complicated by Kasabach-Merritt Syndrome During Pregnancy},
  colorlinks=true,
  linkcolor={blue},
  filecolor={Maroon},
  citecolor={Blue},
  urlcolor={Blue},
  pdfcreator={LaTeX via pandoc}}


\title{Klippel-Trenaunay Syndrome Complicated by Kasabach-Merritt
Syndrome During Pregnancy}
\usepackage{etoolbox}
\makeatletter
\providecommand{\subtitle}[1]{% add subtitle to \maketitle
  \apptocmd{\@title}{\par {\large #1 \par}}{}{}
}
\makeatother
\subtitle{A case report and review of the literature}
\author{}
\date{}
\begin{document}
\maketitle


\section{Authors}\label{authors}

Mourad Elfaham\textsuperscript{1}, Aya Attia\textsuperscript{1}, Mohamed
Waheed\textsuperscript{1}, Mohamed Essam\textsuperscript{1}, Amgad
Gamal\textsuperscript{1}, Salma Saad\textsuperscript{1}, Mariam
Sarhan\textsuperscript{1}, Marwa Mourad\textsuperscript{1}, Mohamed
Hegazy\textsuperscript{1}, Farah Said\textsuperscript{1}, Sarah
Tarek\textsuperscript{1}, Sara Ibrahim Abdelkader\textsuperscript{2*},
Ahmed Azab\textsuperscript{2}, Mohammed Ghanem\textsuperscript{2}, Rahma
Alaa Abdelhafez\textsuperscript{2}, Nour Atif\textsuperscript{2}, Maha
Moemen\textsuperscript{2}, Belgin Ahmed Nagah\textsuperscript{2}, Ashraf
Nabhan\textsuperscript{1,3}

\section{Affiliation}\label{affiliation}

\textsuperscript{1} Faculty of Medicine, Ain Shams University, Cairo,
Egypt.

\textsuperscript{2} Faculty of Medicine, Galala University, Attaka,
Suez, Egypt.

\textsuperscript{3} Faculty of Medicine, MTI University, Cairo, Egypt.
01069663999

* Corresponding author: Sara Ibrahim Abdelkader. Faculty of Medicine,
Galala University, Attaka, Suez, Egypt. Email:
\href{mailto:sara.elzeftawy@gu.edu.eg}{\nolinkurl{sara.elzeftawy@gu.edu.eg}},
ORCID: \href{https://orcid.org/0009-0002-5877-4722}{0009-0002-5877-4722}

\textbf{Word count:} 2500

\section{Abstract}\label{abstract}

The co-occurrence of Klippel--Trenaunay syndrome and Kasabach-Merritt
syndrome in pregnancy is extremely rare and represents a critical
obstetric emergency with significant maternal and fetal risks, demanding
early recognition and immediate multidisciplinary intervention.

We report the case of a 28-year-old pregnant woman with known
Klippel--Trenaunay syndrome who developed Kasabach--Merritt phenomenon
during pregnancy. The fetus suffered severe growth restriction and
absent end-diastolic flow on umbilical artery Doppler. A cesarean
section was performed at 35 weeks of gestation.

This case highlights the complexity of managing pregnancy in women with
Klippel--Trenaunay syndrome complicated by Kasabach Merritt syndrome and
emphasizes the importance of multidisciplinary surveillance and timely
delivery.

\textbf{Keywords:} Klippel--Trenaunay syndrome; angioosteohypertrophy
syndrome; Kasabach Merritt syndrome; fetal growth restriction; high-risk
pregnancy

\section{Introduction}\label{introduction}

Klippel-Trenaunay syndrome (KTS) is a rare congenital disorder affecting
an estimated 2--5 per 100,000 individuals. Its features include a triad
of capillary malformations, venous and lymphatic anomalies, and limb
overgrowth{[}1,2{]}. It was first described in 1900 as Naevus vasculosus
osteohypertrophicus by two French physicians, Maurice Klippel and Paul
Trenaunay. It was previously considered similar to the less common
Parkes Weber syndrome, distinguished from KTS by the presence of
arteriovenous malformations. {[}3,4{]}.

KTS hallmark features predominantly affect one lower limb (85\%) causing
a discrepancy in limb length, although bilateral involvement may occur
in a minority of cases (12.5\%). However, it is a generalized disorder
and can involve multiple body regions including the pelvis, abdomen,
head and neck, back, pelvis, and urinary bladder. Capillary
malformations are the most common presentation (98\%), followed by
varicosities or venous malformations (72\%), and limb hypertrophy (67\%)
{[}5{]}.

Pregnancy poses considerable challenges in patients with KTS, as
hormonal and hemodynamic changes exacerbate underlying vascular
abnormalities increasing the risk of thromboembolism, hemorrhage, and
coagulopathy. The emergence of coagulopathy, in particular, may indicate
the development of Kasabach-Merritt syndrome (KMS). {[}6{]}

Kasabach-Merritt syndrome (KMS), first described in 1940, is a rare and
potentially fatal condition defined by the triad of thrombocytopenia,
microangiopathic hemolytic anemia, and consumptive coagulopathy
occurring in association with vascular abnormalities {[}7{]}. In the
setting of KTS, pregnancy and labor may precipitate or exacerbate KMS;
the concurrent occurrence of both conditions during the peripartum
period constitutes a critical obstetric emergency carrying a reported
mortality rate of 30--40\% {[}8{]}.

To date, only three cases of KTS complicated by KMS in pregnancy have
been reported in the literature {[}9--11{]}. We hereby describe the
clinical course and successful multidisciplinary management of a
pregnancy in a woman with KTS complicated by fetal growth restriction
and peripartum acute KMS. Through this case, we aim to describe the
clinical challenges inherent to this rare case, outline effective
management strategies, and emphasize the critical role of early
recognition and interdisciplinary collaboration in achieving favorable
maternal and neonatal outcomes.

\section{Case Presentation}\label{case-presentation}

A 28-year-old multigravida presented at 30+1 weeks' gestation with
clinically significant hematuria and was admitted on December 16th to
Ain Shams University Hospital, Egypt. Her medical history was notable
for Klippel--Trenaunay syndrome (KTS), diagnosed at two years of age.
She had an extensive KTS phenotype, with disease involvement
encompassing both lower limbs, the back, pelvis, urinary bladder, and
breast tissue --- a distribution reflecting a multisystem vascular
dysplasia. Her surgical history included varicose vein surgery seven
years prior and orthopedic fixation with plates and screws approximately
18 months before this admission. Her past obstetric history was
significant. Her first pregnancy, seven years prior, ended in neonatal
death. Her second pregnancy was complicated by severe preeclampsia,
postpartum wound hematoma, and sepsis necessitating ICU admission. She
also had a history of one spontaneous miscarriage In the current
pregnancy, she received two doses of antenatal corticosteroids on
December 14th and 15th.

On admission, vital signs were stable with a blood pressure 120/80 mmHg,
temperature 36.2°C, and heart rate 68 bpm. Physical examination revealed
port-wine stains over the back and gluteal region, with discrepancy in
lower limb length bilaterally. Blood typing identified B-positive blood
group, and viral serology was negative for hepatitis B, hepatitis C, and
HIV. Her complete blood count (CBC) demonstrated microcytic hypochromic
anemia (hemoglobin 7.0 g/dL, hematocrit 23.0\%) and thrombocytopenia
(128,000/µL). She received six units of packed red blood cells (pRBCs)
to correct her declining hemoglobin. Obstetric ultrasonography revealed
an estimated fetal weight below the 10th percentile, consistent with
fetal growth restriction (FGR). There was an absent end-diastolic flow
on umbilical artery Doppler.

Urological consultation was obtained for management of severe hematuria,
with initial measures comprising intravenous tranexamic acid and
insertion of a three-way indwelling silicone urinary catheter maintained
for one week.

On December 18th, pelviabdominal MRI demonstrated a viable single
intrauterine pregnancy in cephalic presentation with unremarkable
placental location and attachment. Compression of the left ovarian and
left common iliac veins by the gravid uterus was identified, with
associated venous dilatation and slowing of flow. Extensive dilated,
tortuous tubular lesions consistent with varicosities were observed
throughout the pelvic, perivesical, perineal, and gluteal regions ---
more prominent on the left --- exhibiting a classic ``bag of worms''
appearance. Similar lesions were identified in the subcutaneous tissue
of both breasts. The urinary bladder contained blood. The liver was
mildly enlarged with vascular compression, while the gallbladder,
spleen, and kidneys were unremarkable. Minimal bilateral pleural
effusion was also noted.

On December 21st, a pelviabdominal ultrasonography, color Doppler of
bilateral lower limb veins, and duplex arterial studies were performed.
Ultrasonography demonstrated hepatosplenomegaly, soft tissue lesions
within the urinary bladder, and bilateral refluxing parametrial venous
channels, with a single viable intrauterine fetus confirmed. Lower limb
vascular studies revealed a normal venous and arterial system with no
evidence of deep vein thrombosis.

Serial CBCs over the four days following transfusion showed slight
improvement in hemoglobin, reaching 9.4 g/dL on December 22nd, while
platelet count remained persistently low, ranging between 82,000 and
104,000/µL. Iron studies revealed elevated unsaturated iron-binding
capacity (UIBC) and reduced transferrin saturation. Lactate
dehydrogenase (LDH) was elevated at 454 IU/L. Urinalysis on the same day
demonstrated hematuria (2+), pyuria, and mild proteinuria (1+);
microscopic examination showed \textgreater100 red blood cells per
high-power field (RBCs/HPF), 40--45 pus cells/HPF, increased epithelial
cells (2+), and uric acid crystals (2+). Coagulation studies on December
27th and 28th revealed mildly elevated prothrombin time (PT), while
serum creatinine remained low.

Contrast-enhanced pelvic MRI on January 1st demonstrated extensive
venous malformations involving the cervix, vagina, anal canal, perineum,
and bilateral labia, with additional lesions in the gluteal region,
thighs, anterior abdominal wall, and intravesically.

The patient was monitored daily with regular assessment of vital
parameters, fetal kick counts, vaginal bleeding, uterine contractions,
and signs of membrane rupture. She received one unit of pRBCs on January
2nd and a further unit on January 5th.

At 35 weeks' gestation, a multidisciplinary team comprising obstetrics,
hematology, anesthesia, and neonatology convened and recommended
delivery by emergency cesarean section in view of deteriorating fetal
Doppler parameters and cumulative maternal risks. On January 18th, after
preoperative preparation with 2gm ceftriaxone, cesarean delivery was
performed under general anesthesia via a vertical midline abdominal
incision. Intraoperatively, significant varicosities were encountered
over the anterior abdominal wall, lower uterine segment, and the
uterovesical interface. A lower uterine segment incision was made for
delivery of the fetus and placenta; the uterine wall was subsequently
closed in two layers. The placenta was delivered complete, the perineum
was intact, and hemostasis was secured. No intraperitoneal drains were
placed; a subcutaneous drain was inserted and maintained for seven days.
Skin closure was achieved with interrupted mattress sutures. One packed
RBCs, one FFP, 12 units pf platelets and 2 gm tranexamic acid were
infused intraoperatively to control oozing.

A female neonate weighing 1,900 g was delivered with Apgar scores of 7
and 9 at one and five minutes, respectively. Given the low birth weight,
routine vaccination was deferred. The newborn did not require
respiratory support.

Postoperatively, the patient was admitted to the ICU following
identification of hemoperitoneum. Her pharmacological regimen included
low-molecular-weight heparin (LMWH), tranexamic acid, ceftriaxone,
pantoprazole, intravenous paracetamol, and a charcoal-based herbal
preparation.

On postoperative day one, mild hematuria persisted and laboratory
investigations revealed declining hemoglobin and platelet counts
alongside markedly elevated LDH, alkaline phosphatase, creatine kinase,
and bilirubin. LMWH was withheld, and the patient was commenced on pRBCs
and fresh frozen plasma (FFP) in a 1:1 ratio, receiving a total of ten
units of each over seven days. Additional laboratory results are
presented in \textbf{?@tbl-2}

On postoperative day two, hemoglobin and platelet count deteriorated
rapidly, falling from 9.6 g/dL and 218,000/µL to 7.6 g/dL and
104,000/µL, respectively, necessitating ongoing blood product
replacement.

On postoperative day three, transvaginal ultrasonography (TVUS) revealed
marked hemoperitoneum with a localized lower uterine segment hematoma
measuring approximately 12 cm in maximum diameter. A decision to
continue conservative management was made by the attending consultant
and by the head of the unit. The pharmacological regimen was revised to
include LMWH, tranexamic acid, dexamethasone, magnesium sulfate,
cefepime, omeprazole, clindamycin, and intravenous paracetamol.

On postoperative day four, coagulation studies revealed
hypofibrinogenemia (fibrinogen 1.4 g/L) and markedly elevated D-dimer
(35.2 µg/mL FEU), findings consistent with a diagnosis of
Kasabach-Merritt syndrome (KMS).

Urinalysis on day six demonstrated persistent hematuria (2+), mild
pyuria, and mild proteinuria (1+), with microscopic examination showing
\textgreater100 RBCs/HPF, 13--15 pus cells/HPF, and increased epithelial
cells (1+). By day seven, hemoglobin had stabilized at 10.1 g/dL and
remained stable over the subsequent three days, though thrombocytopenia
persisted at 64,000/µL.

Following clinical stabilization, the patient was discharged on
postoperative day ten. At her outpatient review on day nineteen, she was
in satisfactory general condition with a soft, lax abdomen; skin sutures
were removed and the abdominal wound demonstrated optimal healing.

\textbf{?@tbl-1} shows the full clinical timeline in chronological
order.

\begin{longtable}[]{@{}
  >{\raggedright\arraybackslash}p{(\linewidth - 2\tabcolsep) * \real{0.2791}}
  >{\raggedright\arraybackslash}p{(\linewidth - 2\tabcolsep) * \real{0.7209}}@{}}
\caption{Patient timeline}\tabularnewline
\toprule\noalign{}
\begin{minipage}[b]{\linewidth}\raggedright
Date
\end{minipage} & \begin{minipage}[b]{\linewidth}\raggedright
Timeline
\end{minipage} \\
\midrule\noalign{}
\endfirsthead
\toprule\noalign{}
\begin{minipage}[b]{\linewidth}\raggedright
Date
\end{minipage} & \begin{minipage}[b]{\linewidth}\raggedright
Timeline
\end{minipage} \\
\midrule\noalign{}
\endhead
\bottomrule\noalign{}
\endlastfoot
30 +1 weeks' gestation (Dec 16, 2025) & Admission of a case of KTS,
presenting with hematuria; anemia and thrombocytopenia; fetus
\textless10th percentile, received tranexamic acid and 6 units PRBCs. \\
30+1 to 35 weeks' gestation (Dec 18, 2025 -- Jan 17, 2026) &
\begin{minipage}[t]{\linewidth}\raggedright
MRI: extensive venous malformations (pelvis, thigh, abdomen).\\
Close maternal/fetal monitoring.\\
Intermittent transfusions slight hemoglobin improvement; fluctuating
platelets; prolonged PT.\strut
\end{minipage} \\
35 weeks (Jan 18, 2026) & Emergency cesarean section for fetal
compromise → female neonate (1900 g). Mother admitted to ICU,
hemoperitoneum managed conservatively. \\
Postoperative course & \textbf{Day 1:} Transfusion support (total 10
units PRBC + FFP over 7 days)

\textbf{Day 2:} Hb 9.6→7.6 g/dL; platelets 218×10³→104×10³/µL

\textbf{Day 3:} Conservative management, LMWH resumed, tranexamic acid
initiated

\textbf{Day 4:} Hypofibrinogenemia (1.4 g/L), elevated D-dimer (35.2
µg/mL) consistent with Kasabach--Merritt syndrome

\textbf{Day 7:} Hb stabilized (10.1 g/dL); persistent thrombocytopenia
(64×10³/µL)

\textbf{Day 10:} Discharged in stable condition

\textbf{Day 19:} Outpatient follow-up: fair general condition; sutures
removed \\
\end{longtable}

\begin{longtable}[]{@{}
  >{\raggedright\arraybackslash}p{(\linewidth - 16\tabcolsep) * \real{0.1733}}
  >{\raggedright\arraybackslash}p{(\linewidth - 16\tabcolsep) * \real{0.2200}}
  >{\raggedright\arraybackslash}p{(\linewidth - 16\tabcolsep) * \real{0.0933}}
  >{\raggedright\arraybackslash}p{(\linewidth - 16\tabcolsep) * \real{0.0933}}
  >{\raggedright\arraybackslash}p{(\linewidth - 16\tabcolsep) * \real{0.0667}}
  >{\raggedright\arraybackslash}p{(\linewidth - 16\tabcolsep) * \real{0.0667}}
  >{\raggedright\arraybackslash}p{(\linewidth - 16\tabcolsep) * \real{0.0600}}
  >{\raggedright\arraybackslash}p{(\linewidth - 16\tabcolsep) * \real{0.0600}}
  >{\raggedright\arraybackslash}p{(\linewidth - 16\tabcolsep) * \real{0.1200}}@{}}
\caption{Laboratory results throughout the clinical
course}\tabularnewline
\toprule\noalign{}
\begin{minipage}[b]{\linewidth}\raggedright
Parameters/Timing
\end{minipage} & \begin{minipage}[b]{\linewidth}\raggedright
December 16 (Day of admission)
\end{minipage} & \begin{minipage}[b]{\linewidth}\raggedright
December 22
\end{minipage} & \begin{minipage}[b]{\linewidth}\raggedright
December 28
\end{minipage} & \begin{minipage}[b]{\linewidth}\raggedright
POD 1
\end{minipage} & \begin{minipage}[b]{\linewidth}\raggedright
POD 2
\end{minipage} & \begin{minipage}[b]{\linewidth}\raggedright
POD 4
\end{minipage} & \begin{minipage}[b]{\linewidth}\raggedright
POD 6
\end{minipage} & \begin{minipage}[b]{\linewidth}\raggedright
Reference range
\end{minipage} \\
\midrule\noalign{}
\endfirsthead
\toprule\noalign{}
\begin{minipage}[b]{\linewidth}\raggedright
Parameters/Timing
\end{minipage} & \begin{minipage}[b]{\linewidth}\raggedright
December 16 (Day of admission)
\end{minipage} & \begin{minipage}[b]{\linewidth}\raggedright
December 22
\end{minipage} & \begin{minipage}[b]{\linewidth}\raggedright
December 28
\end{minipage} & \begin{minipage}[b]{\linewidth}\raggedright
POD 1
\end{minipage} & \begin{minipage}[b]{\linewidth}\raggedright
POD 2
\end{minipage} & \begin{minipage}[b]{\linewidth}\raggedright
POD 4
\end{minipage} & \begin{minipage}[b]{\linewidth}\raggedright
POD 6
\end{minipage} & \begin{minipage}[b]{\linewidth}\raggedright
Reference range
\end{minipage} \\
\midrule\noalign{}
\endhead
\bottomrule\noalign{}
\endlastfoot
Hb (g/dL) & 7.0 & 9.4 & 10.2 & 9.6 & 7.6 & 8.6 & 10.0 & 12.0-15.0 \\
Platelets (per µL) & 128,000 & 97,000 & 129,000 & 218,000 & 104,000 &
75,000 & 64,000 & 150,000--410,000 \\
PT (seconds) & - & - & 12.4 & 40.70 & - & - & 17.7 & 11.7--15.3 \\
Total bilirubin (mg/dL) & - & - & - & 2.8 & 1.8 & - & 1.4 & 1.8-3.5 \\
ALP (IU/L) & - & - & - & 135 & 115 & - & 108 & 34--104 \\
LDH (IU/L) & - & 454 & - & 424 & 357 & - & 449 & 140--271 \\
\end{longtable}

POD: Post operative day, Hb: Hemoglobin, HTC: Hematocrit, PT:
Prothrombin time, ALP: Alkaline phosphatase, LDH : Lactate
dehydrogenase.

\section{Review of literature}\label{review-of-literature}

A systematic search of MEDLINE, Embase, and Scopus identified only three
reported cases of KTS complicated by peripartum KMS. KTS was diagnosed
prior to pregnancy in two cases {[}9,10{]}, whereas in the third it had
been misdiagnosed as acromegaly and was not recognized until the final
trimester {[}11{]}. All cases presented with classic KTS features ---
multiple varicosities, limb hypertrophy, and hemangiomas --- with one
additionally demonstrating port-wine stains {[}11{]}. Despite treatment,
postpartum laboratory parameters failed to normalize, indicating
refractory coagulopathy consistent with KMS.

All three cases had marked vulvar varicosities necessitating cesarean
delivery, with intraoperative or postoperative hemorrhage requiring
multiple transfusions; two required massive transfusion protocols
{[}9,11{]}. Management followed a multidisciplinary approach in all
cases, incorporating anticoagulation and blood product support. Two
cases received corticosteroids and antifibrinolytics {[}9,10{]}. One
patient developed life-threatening hemorrhage requiring postpartum
hysterectomy, complicated further by DIC, sepsis, peritonitis, and
hemoperitoneum necessitating emergency laparotomy {[}11{]}. All patients
required prolonged postoperative hospitalization. None of the neonates
demonstrated clinical or laboratory evidence of KTS.

Summary of the three cases is shown in \textbf{?@tbl-3}

\begin{longtable}[]{@{}
  >{\raggedright\arraybackslash}p{(\linewidth - 10\tabcolsep) * \real{0.0723}}
  >{\raggedright\arraybackslash}p{(\linewidth - 10\tabcolsep) * \real{0.1064}}
  >{\raggedright\arraybackslash}p{(\linewidth - 10\tabcolsep) * \real{0.3043}}
  >{\raggedright\arraybackslash}p{(\linewidth - 10\tabcolsep) * \real{0.2319}}
  >{\raggedright\arraybackslash}p{(\linewidth - 10\tabcolsep) * \real{0.1979}}
  >{\raggedright\arraybackslash}p{(\linewidth - 10\tabcolsep) * \real{0.0787}}@{}}
\caption{Literature review}\tabularnewline
\toprule\noalign{}
\begin{minipage}[b]{\linewidth}\raggedright
\textbf{Author (Year)}
\end{minipage} & \begin{minipage}[b]{\linewidth}\raggedright
Pregnancy
\end{minipage} & \begin{minipage}[b]{\linewidth}\raggedright
\textbf{Key findings}
\end{minipage} & \begin{minipage}[b]{\linewidth}\raggedright
\textbf{KMS}
\end{minipage} & \begin{minipage}[b]{\linewidth}\raggedright
\textbf{Management}
\end{minipage} & \begin{minipage}[b]{\linewidth}\raggedright
\textbf{Outcome}
\end{minipage} \\
\midrule\noalign{}
\endfirsthead
\toprule\noalign{}
\begin{minipage}[b]{\linewidth}\raggedright
\textbf{Author (Year)}
\end{minipage} & \begin{minipage}[b]{\linewidth}\raggedright
Pregnancy
\end{minipage} & \begin{minipage}[b]{\linewidth}\raggedright
\textbf{Key findings}
\end{minipage} & \begin{minipage}[b]{\linewidth}\raggedright
\textbf{KMS}
\end{minipage} & \begin{minipage}[b]{\linewidth}\raggedright
\textbf{Management}
\end{minipage} & \begin{minipage}[b]{\linewidth}\raggedright
\textbf{Outcome}
\end{minipage} \\
\midrule\noalign{}
\endhead
\bottomrule\noalign{}
\endlastfoot
\textbf{Neubert (1995)} {[}9{]} & 22 years-old primigravida, 37 weeks'
gestation & known KTS, hemangiomas, limb hypertrophy, massive
vulvovaginal varicosities & Post-op KMS: thrombocytopenia, decreased
fibrinogen, prolonged PT/PTT, increased fibrin split products & CS at 39
wks; hematoma evacuation; ICU; transfusion; heparin + aminocaproic acid
& Discharged day 20; healthy neonate \\
\textbf{Kanmani (2021)} {[}10{]} & 24 years-old, primigravida, 29 weeks'
gestation & known KTS, bleeding gums, Limb hypertrophy, varicosities
(vulvar, lower limb) & Antenatal KMS: severe thrombocytopenia (limited
laboratory investigation reported) & Steroids + platelet transfusion; CS
at 35 wks; post-op anticoagulation and corticosteroids & Discharged day
20; healthy neonate \\
\textbf{Zhang (2019)} {[}11{]} & 26 years old.

Postpartum. & KTS diagnosed in the third trimester of pregnancy.

Admitted on day 10 postpartum with severe postpartum hemorrahge,
hemodynamic instability, ecchymoses, massive limb hypertrophy,
varicosities & Post-op KMS: thrombocytopenia, ↓ fibrinogen, ↑ PT/APTT,
increased fibrin split products, increased D-dimer & CS at 39 wks;
hysterectomy; embolization; laparotomy; anticoagulation & Discharged day
42; healthy neonate \\
\end{longtable}

\section{Discussion}\label{discussion}

The occurrence of Klippel-Trenaunay syndrome complicated by
Kasabach-Merritt syndrome during pregnancy is exceedingly rare, with
only three such cases reported in the literature prior to this report
{[}9--11{]}. The present case thus represents the fourth documented
instance of this life-threatening combination, and the first to describe
its successful management in the context of concomitant fetal growth
restriction and peripartum acute KMS.

\textbf{Diagnosis and baseline complexity}

Our patient presented with an unusually extensive KTS phenotype, with
disease involvement encompassing both lower limbs, the back, pelvis,
urinary bladder, and breast tissue --- a distribution reflecting the
multisystem vascular dysplasia characteristic of KTS. Consistent with
previously reported cases {[}9--11{]}, she exhibited the hallmark triad
of varicosities, limb hypertrophy, and hemangiomas, with port-wine
stains additionally present over the gluteal region. Her obstetric
history further compounded the clinical complexity: two prior cesarean
sections, a neonatal death, a previous pregnancy complicated by severe
preeclampsia and ICU admission, and a spontaneous miscarriage ---
underscoring the cumulative maternal morbidity associated with KTS
across successive pregnancies.

Clinically significant hematuria at 30+1 weeks' gestation was the
sentinel presentation leading to admission, a manifestation attributable
to intravesical venous malformations confirmed on both MRI and
contrast-enhanced pelvic imaging. Notably, pelviabdominal MRI revealed
extensive pelvic varicosities with a characteristic ``bag of worms''
appearance, compression of the left ovarian and common iliac veins by
the gravid uterus, hepatosplenomegaly, and bilateral pleural effusion
--- findings collectively reflecting the hemodynamic burden imposed by
the gravid state on an already compromised vascular system. This aligns
with prior observations that pregnancy-associated hormonal and
hemodynamic changes exacerbate the vascular pathology inherent to KTS,
amplifying risks of thromboembolism and hemorrhage {[}12{]}.

\textbf{Antenatal management}

Initial antenatal management was guided by two principal concerns:
controlling hematuria and optimizing fetal and maternal status ahead of
inevitable preterm delivery. Hematuria was managed conservatively with
intravenous tranexamic acid and continuous bladder irrigation via a
three-way urinary catheter, achieving partial resolution without
surgical intervention. Serial CBCs demonstrated persistent
thrombocytopenia throughout the antenatal period, with platelet counts
ranging between 82,000 and 128,000/µL despite transfusion support,
accompanied by elevated LDH --- early biochemical signals of evolving
consumptive coagulopathy. These antenatal laboratory aberrations, in
retrospect, likely reflected subclinical KMS activity preceding its
overt peripartum manifestation.

Fetal surveillance revealed IUGR with absent end-diastolic umbilical
artery flow, indicating significant uteroplacental insufficiency. The
administration of two doses of antenatal corticosteroids at 30 weeks
represented an important adjunct in preparing for anticipated preterm
delivery. The decision to deliver at 35 weeks was reached through
multidisciplinary consensus, balancing the risk of ongoing fetal
compromise against the substantial maternal surgical risks posed by the
extensive pelvic varicosities. This mirrors the management philosophy
reported in all previously published cases, where cesarean delivery was
ultimately necessitated by vulval or pelvic varicosities precluding safe
vaginal birth {[}9--11{]}.

\textbf{Intraoperative findings and delivery}

Cesarean section was performed under general endotracheal anesthesia,
with the choice of a vertical midline incision allowing adequate
visualization and surgical access in the setting of extensive anterior
abdominal wall and lower uterine segment varicosities. Intraoperatively,
significant varicosities were encountered at the uterovesical interface
--- a finding consistent with the MRI-demonstrated pelvic venous
malformations and a recognized source of hemorrhagic risk in KTS
patients undergoing cesarean delivery. Hemostasis was secured without
recourse to uterine artery ligation, balloon tamponade, or hysterectomy,
distinguishing the intraoperative course of this case favorably from
that of a previously published case report {[}11{]}, in which
life-threatening hemorrhage required postpartum hysterectomy along with
emergency laparotomy for DIC, sepsis, peritonitis, and hemoperitoneum.

A female neonate weighing 1,900 g was delivered with reassuring Apgar
scores of 7 and 9 at one and five minutes, respectively. Consistent with
all three previously reported cases {[}9--11{]}, the neonate
demonstrated no clinical features of KTS, supporting the current
understanding that KTS does not follow a straightforward vertical
transmission pattern.

\textbf{Peripartum KMS and postoperative course}

The postoperative course was complicated by hemoperitoneum requiring ICU
admission, with a rapidly deteriorating hemoglobin and platelet count on
days one and two, accompanied by markedly elevated LDH, alkaline
phosphatase, creatine kinase, and bilirubin. On postoperative day three,
TVUS confirmed a 12-cm lower uterine segment hematoma. The definitive
diagnosis of KMS was established on day four, supported by
hypofibrinogenemia (1.4 g/L) and markedly elevated D-dimer (35.2 µg/mL
FEU) --- laboratory findings fulfilling the diagnostic criteria of
consumptive coagulopathy in the context of underlying vascular
abnormalities. {[}7{]}

Management of peripartum KMS in this case required careful navigation of
the competing hemostatic risks inherent to the syndrome:
thrombocytopenia and consumptive coagulopathy predispose to hemorrhage,
while the vascular malformations underlying KMS simultaneously confer
thrombotic risk. LMWH was initially withheld in response to hemorrhagic
deterioration, then judiciously resumed alongside tranexamic acid once
the coagulopathic picture became clearer. This
anticoagulation-antifibrinolytic combination reflects the dual-pathway
pathophysiology of KMS and is consistent with the approach described in
two of the three previously reported cases {[}9,10{]}. Corticosteroid
therapy, in the form of dexamethasone, was incorporated given its
proposed role in reducing vascular lesion activity and platelet
trapping. The patient ultimately received ten units each of pRBCs and
FFP in a 1:1 ratio over seven days --- constituting a massive
transfusion --- alongside calcium, magnesium sulfate, and broad-spectrum
antibiotics for infective prophylaxis.

Hemoglobin stabilized at 10.1 g/dL by day seven, though thrombocytopenia
persisted at 64,000/µL at discharge --- a pattern consistent with the
refractory coagulopathy observed in all previously reported cases, where
postpartum laboratory parameters failed to normalize despite active
treatment {[}9--11{]}. The patient was discharged on day ten and
demonstrated satisfactory wound healing at her day-nineteen outpatient
review, representing a favorable outcome relative to the severe
postoperative morbidity --- including sepsis, peritonitis, DIC, and
hysterectomy --- reported in the most complicated previously published
case. {[}11{]}

\textbf{Multidisciplinary care and broader implications}

The successful outcome in this case was underpinned by early
multidisciplinary involvement across obstetrics, hematology, anesthesia,
neonatology, urology, and radiology --- a consensus echoed across all
published cases and reflective of the inherent complexity of managing
concurrent KTS and KMS in pregnancy. Early recognition of KMS, even in
its subclinical antenatal form, appears critical: the persistent
antenatal thrombocytopenia and elevated LDH in our patient, though not
initially diagnostic, warranted heightened vigilance and likely informed
the decision to deliver at 35 rather than a later gestational age.

This case further highlights the importance of preconception counseling
in women with KTS. The published cross-sectional data {[}12{]}
documenting significantly elevated risks of DVT, pulmonary embolism, and
severe postpartum hemorrhage in KTS patients provides epidemiological
context for the individual complexity illustrated here. Our patient's
cumulative obstetric history reflects the substantial reproductive
morbidity that KTS may confer across successive pregnancies and
underscores the need for individualized, specialist-led preconception
assessment in this population.

\section{Conclusion}\label{conclusion}

Pregnancy in women with Klippel--Trenaunay syndrome complicated by
Kasabach--Merritt syndrome is extremely rare and poses significant
diagnostic and therapeutic challenges. Early suspicion and diagnosis
with intensive antenatal surveillance and multidisciplinary management
are crucial to navigate the complex balance between thrombosis and
bleeding and to optimize maternal and perinatal outcomes in such
high-risk pregnancy and childbirth.

\section{Abbreviation}\label{abbreviation}

Klippel-Trenaunay syndrome (KTS)

Kasabach Meritt Syndrome (KMS)

Magnetic resonance imaging (MRI)

Ultrasonography (US)

Transvaginal ultrasound (TVUS)

Intensive care unit (ICU)

Low-molecular-weight heparin (LMWH)

Lactate dehydrogenase (LDH)

Complete blood picture (CBC)

Disseminated Intravascular Coagulation (DIC)

Multidisciplinary team (MDT)

\section{Declarations}\label{declarations}

\subsection{Ethics approval and consent to
participate}\label{ethics-approval-and-consent-to-participate}

Not applicable.

\subsection{Consent for publication}\label{consent-for-publication}

A written consent for publication has been obtained from the patient.

\subsection{Availability of data and
materials}\label{availability-of-data-and-materials}

The data and materials supporting the conclusions of this article are
available as a supplementary materials.

\subsection{Competing interests}\label{competing-interests}

The authors declare that they have no competing interests.

\subsection{Funding}\label{funding}

This research received no specific grant from any funding agency in the
public, commercial or not-for-profit sectors.

\subsection{Authors' contributions}\label{authors-contributions}

SIA, ME, AFN contributed to writing the first draft of the manuscript.
All the authors revised the manuscript critically for important
intellectual content. All authors approved the final version of the
manuscript.

\subsection{Acknowledgements}\label{acknowledgements}

We sincerely thank the healthcare providers at Ain Shams University
Hospital of Obstetrics and Gynecology for their dedication and
commitment to patient care in the operating room. We deeply appreciate
their contributions and the valuable role they play in ensuring
high-quality surgical care.

\section*{References}\label{references}
\addcontentsline{toc}{section}{References}

\protect\phantomsection\label{refs}
\begin{CSLReferences}{0}{1}
\bibitem[\citeproctext]{ref-ISSVA2018}
1. International Society for the Study of Vascular Anomalies. {ISSVA
Classification for Vascular Anomalies}.
\url{https://www.issva.org/UserFiles/file/ISSVA-Classification-2018.pdf};
2018.

\bibitem[\citeproctext]{ref-Pavone2023}
2. Pavone P, Marino L, Cacciaguerra G, Di Nora A, Parano E, Musumeci G,
et al. \href{https://doi.org/10.3390/children10081421}{Klippel-trenaunay
syndrome, segmental/focal overgrowth malformations: A review}. Children
(Basel). 2023;10:1421.

\bibitem[\citeproctext]{ref-mckusick1963}
3. McKusick VA. Frederick parkes weber{\textemdash}1863-1962. JAMA
{[}Internet{]}. 1963;183. Available from:
\url{http://dx.doi.org/10.1001/jama.1963.63700010018015}

\bibitem[\citeproctext]{ref-sikova1999}
4. Sikova DVD, Pavlova LjT, Laskoska MTV, Nikolovska ST, Biljanovska NC.
Naevus varicosus osteohypertrophicus. Springer US; 1999. p. 535--40.
Available from: \url{http://dx.doi.org/10.1007/978-1-4615-4857-7_80}

\bibitem[\citeproctext]{ref-Jacob1998}
5. Jacob AG, Driscoll DJ, Shaughnessy WJ, Stanson AW, Clay RP, Gloviczki
P. Klippel-trénaunay syndrome: Spectrum and management. Mayo Clinic
Proceedings {[}Internet{]}. 1998;73:28--36. Available from:
\url{http://dx.doi.org/10.1016/S0025-6196(11)63615-X}

\bibitem[\citeproctext]{ref-Blei2024}
6. Blei F. {Klippel-Trenaunay Syndrome: Clinical Manifestations,
Diagnosis, and Management}. UpToDate.
\url{https://www.uptodate.com/contents/klippel-trenaunay-syndrome-clinical-manifestations-diagnosis-and-management};
2024.

\bibitem[\citeproctext]{ref-Mahajan2017}
7. Mahajan P, Margolin J, Iacobas I. Kasabach-merritt phenomenon:
Classic presentation and management options. Clinical Medicine Insights:
Blood Disorders {[}Internet{]}. 2017;10:1179545X1769984. Available from:
\url{http://dx.doi.org/10.1177/1179545X17699849}

\bibitem[\citeproctext]{ref-Tian2026}
8. Tian Y, Yuan Y, Wei L, Xu Z, Li L, Li L.
\href{https://doi.org/10.1002/ped4.70041}{Kaposiform
hemangioendothelioma: Diagnosis and treatment}. Pediatric Investigation.
2026;

\bibitem[\citeproctext]{ref-Neubert1995}
9. Neubert GA, Golden MA, Rose NC.
\href{https://doi.org/10.1016/0029-7844(94)00349-I}{Kasabach-merritt
coagulopathy complicating klippel-trenaunay-weber syndrome in
pregnancy}. Obstetrics \& Gynecology. 1995;85:831--3.

\bibitem[\citeproctext]{ref-Kanmani2021}
10. Kanmani K, Meena M, Narmadha D, A. GV.
\href{https://doi.org/10.18203/2320-1770.ijrcog20211511}{Klippel
trenaunay syndrome--an obstetric challenge}. International Journal of
Reproduction, Contraception, Obstetrics and Gynecology. 2021;10:2081--4.

\bibitem[\citeproctext]{ref-zhang2019}
11. Zhang J, Wang K, Mei J. Late puerperal hemorrhage of a patient with
klippel{\textendash}trenaunay syndrome. Medicine {[}Internet{]}.
2019;98:e18378. Available from:
\url{http://dx.doi.org/10.1097/MD.0000000000018378}

\bibitem[\citeproctext]{ref-silvacorreia2024}
12. Silva Correia IF, Hussain M, Johnson J-A. Obstetric management for
pregnant women with klippel{\textendash}trenaunay syndrome: A UK case
report and review of the literature. International Journal of Gynecology
\& Obstetrics {[}Internet{]}. 2024;168:484--6. Available from:
\url{http://dx.doi.org/10.1002/ijgo.15889}

\end{CSLReferences}




\end{document}
